\documentclass[UTF8]{ctexart}
\usepackage{xeCJK}
\usepackage{geometry}
\geometry{left=2cm,right=2cm,top=2.5cm,bottom=2.5cm}
\setlength{\headheight}{12.7pt}

\usepackage{titlesec}
\titleformat{\section}
  {\normalfont\large\bfseries}
  {\thesection}
  {1em}
  {}
\titlespacing*{\section}
  {0pt}
  {3.5ex plus 1ex minus .2ex}
  {2.3ex plus .2ex}

\usepackage{amsmath}
\usepackage{amssymb}
\usepackage{amsthm}
\usepackage{enumitem}
\setlist[enumerate]{itemsep=0pt,topsep=0pt,parsep=0pt,leftmargin=0pt,itemindent=2.5em}
\setlist[itemize]{itemsep=0pt,topsep=0pt,parsep=0pt,leftmargin=0pt,itemindent=2.5em}
\usepackage{fancyhdr}
\pagestyle{fancy}
\fancyhf{}
\usepackage{graphicx}
\usepackage{hyperref}
\usepackage{indentfirst}
\usepackage{lmodern}


\begin{document}
\setlength{\abovedisplayskip}{1pt}
\setlength{\belowdisplayskip}{1pt}

\section{部分正确性公式}

\hypertarget{sec:定义1.1}{}
\noindent\textbf{定义1.1 (部分正确性公式)}

对任意的断言(即公式)$P,Q$和while程序$S$, 称
\[
    \{P\}\,S\,\{Q\}
\]
为一个\textbf{部分正确性公式}.

\vspace{10pt}

\hypertarget{sec:定义1.2}{}
\noindent\textbf{定义1.2 (部分正确性公式的语义)}

给定断言$P,Q$和程序$S$, 如果$\mathcal{M}[\![S]\!](P)\subset Q$, 即
\[
    \forall\sigma,\tau\in\Sigma:\quad
    \,\sigma\models P\land [\![S]\!](\sigma)=\tau\quad\to\quad\tau\models Q\,,
\]
就记$\models\{P\}\,S\,\{Q\}$.

\vspace{10pt}

\hypertarget{sec:定义1.3}{}
\noindent\textbf{定义1.3 (PW证明系统)}

while程序部分正确性公式的PW(partial, while)证明系统:
\begin{enumerate}
    \item \textbf{skip公理}:
    \[
        \{P\}\,\mathrm{skip}\,\{P\},
    \]
    \item \textbf{assignment公理}(后向形式):
    \[
        \{\,P[u:=t]\,\}\,u:=t\,\{P\},
    \]
    \item \textbf{composition规则}:
    \[
        \frac{\{P\}\,S_1\,\{R\},\,\{R\}\,S_2\,\{Q\}}
        {\{P\}\,S_1;S_2\,\{Q\}},
    \]
    \item \textbf{conditional规则}:
    \[
        \frac{\{P\land B\}\,S_1\,\{Q\},\,\{P\land\neg\,B\}\,S_2\,\{Q\}}
        {\{P\}\,\mathbf{if}\,\,B\,\,\mathbf{then}\,\,S_1\,\,\mathbf{else}
        \,\,S_2\,\,\mathbf{fi}\{Q\}},
    \]
    \item \textbf{loop规则}:
    \[
        \frac{\{P\land B\}\,S\,\{P\}}
        {\{P\}\,\mathbf{while}\,\,B\,\,\mathbf{do}\,\,S\,\,\mathbf{od}\,
        \{P\land\neg\,B\}},
    \]
    \item \textbf{consequence规则}:
    \[
        \frac{P\to P_1,\,\{P_1\}\,S\,\{Q_1\},\,Q_1\to Q}
        {\{P\}\,S\,\{Q\}}.
    \]
\end{enumerate}

\vspace{15pt}

assignment公理有前向形式: 当与$u$同类型的变元$v$不出现在$u,t,P$中时,
\[
    \vdash_{\mathrm{PW}}
    \{P\}\,u:=t\,\{\,\exists v\,(\,u\equiv t[u:=v]\land P[u:=v]\,)\,\}.
\]
它可由后向形式推出: 记$Q=\exists v\,(\,u\equiv t[u:=v]\land P[u:=v]\,)$, 则易证
\begin{enumerate}
    \item $P\to t\equiv t[u:=u]\land P[u:=u]\to
    \exists v\,(\,t\equiv t[u:=v]\land P[u:=v]\,)$,
    \item $Q[u:=t]=\exists v\,(\,t\equiv t[u:=v]\land P[u:=v]\,)$,
\end{enumerate}
所以$P\to Q[u:=t]$且$\vdash_{\mathrm{PW}}\{Q[u:=t]\}\,u:=t\,\{Q\}$,
依consequence规则得$\vdash_{\mathrm{PW}}\{P\}\,u:=t\,\{Q\}$.





\newpage

\section{完全正确性公式}

\hypertarget{sec:定义2.1}{}
\noindent\textbf{定义2.1 (完全正确性公式)}

对任意的断言(即公式)$P,Q$和while程序$S$, 称
\[
    [P]\,S\,[Q]
\]
为一个\textbf{完全正确性公式}.

\vspace{10pt}

\hypertarget{sec:定义2.2}{}
\noindent\textbf{定义2.2 (完全正确性公式的语义)}

给定断言$P,Q$和程序$S$, 如果$\mathcal{M}_{\mathrm{tot}}[\![S]\!](P)\subset Q$, 即
\[
    \forall\sigma\in\Sigma:\quad\,\sigma\models P
    \quad\to\quad[\![S]\!](\sigma)\models Q\,,
\]
就记$\models[P]\,S\,[Q]$.

\vspace{10pt}

\hypertarget{sec:定义2.3}{}
\noindent\textbf{定义2.3 (TW证明系统)}

while程序完全正确性公式的TW(total, while)证明系统:
\begin{enumerate}
    \item \textbf{skip公理}:
    \[
        [P]\,\mathrm{skip}\,[P],
    \]
    \item \textbf{assignment公理}:
    \[
        [\,P[u:=t]\,]\,u:=t\,[P],
    \]
    \item \textbf{composition规则}:
    \[
        \frac{[P]\,S_1\,[R],\,[R]\,S_2\,[Q]}
        {[P]\,S_1;S_2\,[Q]},
    \]
    \item \textbf{conditional规则}:
    \[
        \frac{[P\land B]\,S_1\,[Q],\,[P\land\neg\,B]\,S_2\,[Q]}
        {[P]\,\mathbf{if}\,\,B\,\,\mathbf{then}\,\,S_1\,\,\mathbf{else}
        \,\,S_2\,\,\mathbf{fi}[Q]},
    \]
    \item \textbf{loop规则II}: 当$t$是整型项, $z$是整型变量,
    且$z$不出现在$t,B,P,S$中时,
    \vspace{5pt}
    \[
        \frac{[P\land B]\,S\,[P],\,[P\land B\land t=z]\,S\,[t<z],\,
        P\to t\geqslant 0}
        {[P]\,\mathbf{while}\,\,B\,\,\mathbf{do}\,\,S\,\,\mathbf{od}\,
        [P\land\neg\,B]}\,,
    \]
    \item \textbf{consequence规则}:
    \[
        \frac{P\to P_1,\,[P_1]\,S\,[Q_1],\,Q_1\to Q}
        {[P]\,S\,[Q]}.
    \]
\end{enumerate}

\vspace{15pt}

\hypertarget{sec:定理2.1}{}
\noindent\textbf{定理2.1}

对任意的断言$P,Q$和程序$S$, 如果$\vdash_{\mathrm{TW}}[P]\,S\,[Q]$,
那么$\vdash_{\mathrm{PW}}\{P\}\,S\,\{Q\}$.

\noindent{\kaishu 证明.}

对TW规则归纳证明, 显然. \hfill $\qedsymbol$

\end{document}
